\documentclass[onecolumn,aps,prb,amsmath,amssymb,superscriptaddress,longbibliography]{revtex4-2}

\usepackage{newtxtext}
\usepackage{newtxmath}
\usepackage{microtype}
\usepackage{bm}
\usepackage{mathtools}
\usepackage{xcolor}
\definecolor{PRLBlue}{RGB}{46,48,146}
\usepackage[
    colorlinks,
    citecolor=PRLBlue,
    linkcolor=PRLBlue,
    urlcolor=PRLBlue,
]{hyperref}
\usepackage{orcidlink}

\newcommand{\Tr}{\operatorname{Tr}}
\newcommand{\diag}{\operatorname{diag}}
\newcommand{\ee}{\mathrm{e}}
\newcommand{\ii}{\mathrm{i}}
\newcommand{\trans}{\mathsf{T}}
\newcommand{\mc}[1]{\mathcal{#1}}
\newcommand{\bb}[1]{\bm{#1}}
\newcommand{\ror}[2]{\href{https://ror.org/#1}{#2}}

\allowdisplaybreaks

\begin{document}

\title{Time-Reversal Symmetries for Sign-Problem-Free Quantum Monte Carlo}

\author{Yin-Kai Yu\,\orcidlink{0009-0006-2589-2772}}
\affiliation{Beijing National Laboratory for Condensed Matter Physics\\ \& \ror{05cvf7v30}{Institute of Physics, Chinese Academy of Sciences}, Beijing 100190, China}
\affiliation{\ror{05qbk4x57}{University of Chinese Academy of Sciences}, Beijing 100049, China}

\author{Zhong-Chao Wei\,\orcidlink{0000-0002-1708-4306}}
\email{zcwei@thp.uni-koeln.de}
\affiliation{Qingdao Binhai University, Qingdao 266555, China}

\author{Lei Wang\,\orcidlink{0009-0004-1348-8583}}
\email{wanglei@iphy.ac.cn}
\affiliation{Beijing National Laboratory for Condensed Matter Physics\\ \& \ror{05cvf7v30}{Institute of Physics, Chinese Academy of Sciences}, Beijing 100190, China}
\affiliation{\ror{05qbk4x57}{University of Chinese Academy of Sciences}, Beijing 100049, China}
\affiliation{\ror{020vtf184}{Songshan Lake Materials Laboratory}, Dongguan, Guangdong 523808, China}

\author{Zi-Xiang Li\,\orcidlink{0000-0002-3941-365X}}
\email{zixiangli@iphy.ac.cn}
\affiliation{Beijing National Laboratory for Condensed Matter Physics\\ \& \ror{05cvf7v30}{Institute of Physics, Chinese Academy of Sciences}, Beijing 100190, China}
\affiliation{\ror{05qbk4x57}{University of Chinese Academy of Sciences}, Beijing 100049, China}

\date{\today}

\begin{abstract}
Auxiliary-field quantum Monte Carlo reduces an interacting problem to a statistical sum over products of Gaussian many-body operators.  We classify when mutually anticommuting time-reversal symmetries make every bosonic or fermionic configuration weight nonnegative, including propagators with pairing terms.  Each time reversal squares to \(+I\) or \(-I\), and braces list these signs.  The Majorana class \(\{+,-\}\), containing one symmetry of each square, separates the evolution into complex-conjugate sectors.  The Kramers class \(\{-,-\}\) selects one operator from every Nambu-conjugate pair and imposes Kramers degeneracy.  Both classes protect bosonic and fermionic weights.  The bosonic split-orthogonal class \(\{+,+,+\}\) provides a third mechanism, while the same symmetry set permits negative fermionic weights.  Four or more anticommuting time reversals always generate a Majorana or Kramers pair.  The resulting periodic table unifies time-reversal criteria for both statistics and guides sign-free interacting boson models with anomalous propagation.
\end{abstract}

\maketitle

\tableofcontents

\section{Quadratic evolution and configuration weights}

A Hubbard--Stratonovich decomposition expresses an interacting partition function as an integral or sum over auxiliary-field configurations~\cite{blankenbecler1981prd}.  Once one configuration is fixed, the quantum degrees of freedom evolve through a sequence of quadratic generators.  When the many-body trace converges, the corresponding Monte Carlo weight is
\begin{equation}
    p=\Tr\!\left(\ee^{\hat h_1}\ee^{\hat h_2}\cdots\ee^{\hat h_L}\right),
    \qquad
    \hat h_\ell=\frac12\hat\chi^\trans gA_\ell\hat\chi .
    \label{eq:configuration_weight}
\end{equation}
Here
\begin{equation}
    \hat\chi=
    (\hat a_1,\ldots,\hat a_N,
     \hat a_1^+,\ldots,\hat a_N^+)^\trans
    \label{eq:nambu_spinor}
\end{equation}
contains the annihilation and creation operators.  The same notation describes bosons and fermions when
\begin{equation}
    g_{\rm B}=\begin{pmatrix}0&-I_N\\ I_N&0\end{pmatrix},
    \qquad
    g_{\rm F}=\begin{pmatrix}0&I_N\\ I_N&0\end{pmatrix}.
    \label{eq:g_matrices}
\end{equation}
The condition
\begin{equation}
    gA_\ell+A_\ell^\trans g=0
    \label{eq:quadratic_algebra}
\end{equation}
makes \(gA_\ell\) symmetric for bosons and antisymmetric for fermions, exactly as required by their commutation relations.  The generators \(\hat h_\ell\) may be non-Hermitian after the Hubbard--Stratonovich transformation even when the original interacting Hamiltonian is Hermitian.

The action of the ordered Gaussian product on the Nambu operators is encoded by
\begin{equation}
    M=\ee^{A_1}\ee^{A_2}\cdots\ee^{A_L}.
    \label{eq:propagation_matrix}
\end{equation}
The bosonic Gaussian trace identity~\cite{munn1978ensemble} and its paired-fermion counterpart~\cite{klich2014jstat} are
\begin{equation}
    p_{\rm B}=\left[(-1)^N\det(I_{2N}-M)\right]^{-1/2},
    \qquad
    p_{\rm F}=\left[\det(I_{2N}+M)\right]^{1/2}.
    \label{eq:nambu_trace}
\end{equation}
Their normalization follows already from one mode.  If the two Nambu eigenvalues are \(\ee^x\) and \(\ee^{-x}\), the centered fermion generator gives
\begin{equation}
    \Tr\ee^{x(\hat n-1/2)}
    =\ee^{-x/2}(1+\ee^x),
    \label{eq:fermion_one_mode}
\end{equation}
while the boson generator gives
\begin{equation}
    \Tr\ee^{x(\hat n+1/2)}
    =\frac{\ee^{x/2}}{1-\ee^x}.
    \label{eq:boson_one_mode}
\end{equation}
The determinant inside each square root in Eq.~\eqref{eq:nambu_trace} is the square of the corresponding mode product.  The one-body matrix \(M\) alone has therefore lost one overall sign.  The ordered Fock-space operators in Eq.~\eqref{eq:configuration_weight} fix that sign~\cite{klich2014jstat}.  Starting in the trace-convergent region and varying the same ordered operators continuously also defines the Gaussian weight when an occupation-number sum diverges.  We use \(p\) for this continuation as well as for the ordinary trace with which it agrees.  For bosons, \(\det(I_{2N}-M)=0\) is a pole.  All bosonic positivity statements below apply away from these poles.

We use time reversal in the generalized quantum Monte Carlo sense of an antiunitary symmetry of every auxiliary-field slice.  Let \(\widehat\Theta\) be its action on Fock space and use \(\Theta\) for its anti-linear action on the space of one-body operators,
\begin{equation}
    \Theta(\hat\chi)
    \equiv\widehat\Theta\hat\chi\widehat\Theta^{-1}
    =U\hat\chi,
    \qquad
    \Theta(c\hat d)=\overline c\,\Theta(\hat d).
    \label{eq:time_reversal_definition}
\end{equation}
Here \(c\) is a complex number and \(\hat d\) is any linear combination of the entries of \(\hat\chi\).  The rows of \(U\) record the images of the basis operators in \(\hat\chi\).  We use the compact label \(\Theta=U\mc K\), with Eq.~\eqref{eq:time_reversal_definition} fixing this operator-space convention.  The order of two transformations then follows directly:
\begin{equation}
    (\Theta_a\Theta_b)(\hat\chi)
    =\Theta_a(U_b\hat\chi)
    =\overline U_bU_a\hat\chi.
    \label{eq:TR_composition_order}
\end{equation}
The complex conjugation acts on \(U_b\) because \(\Theta_a\) is anti-linear.  The Nambu exchange matrix
\begin{equation}
    C=\begin{pmatrix}0&I_N\\ I_N&0\end{pmatrix}
    \label{eq:C_matrix}
\end{equation}
encodes Hermitian conjugation between the upper and lower halves of \(\hat\chi\).  A physical time reversal satisfies
\begin{equation}
    U^\trans gU=g,
    \qquad
    C\overline U C=U,
    \qquad
    U^{-1}\overline{A_\ell}U=A_\ell
    \quad\text{for every }\ell .
    \label{eq:physical_TR_conditions}
\end{equation}
The first relation preserves the bosonic commutator or fermionic anticommutator.  The second ensures that the transformed creation operator remains the Hermitian conjugate of the transformed annihilation operator.  The third keeps the fixed-field quadratic generator invariant.  Indeed,
\begin{equation}
    \widehat\Theta\hat h_\ell\widehat\Theta^{-1}
    =\frac12\hat\chi^\trans U^\trans g\overline{A_\ell}U\hat\chi
    =\frac12\hat\chi^\trans gA_\ell\hat\chi
    =\hat h_\ell .
    \label{eq:TR_invariance_derivation}
\end{equation}
The third relation can be rewritten as \(\overline{A_\ell}U=UA_\ell\).  It shows directly that time reversal commutes with the infinitesimal one-body evolution,
\begin{equation}
    \Theta(A_\ell\hat\chi)
    =\overline{A_\ell}U\hat\chi
    =UA_\ell\hat\chi.
    \label{eq:TR_commutes_A}
\end{equation}
Applying \(\Theta\) twice to \(\hat\chi\) gives
\begin{equation}
    \Theta^2(\hat\chi)=\overline U U\hat\chi
    =\pm\hat\chi.
    \label{eq:TR_square}
\end{equation}
We denote these two possibilities by \(+\) and \(-\).  A negative square gives Kramers partners.  If \(\Theta\hat d=c\hat d\), applying \(\Theta\) once more would give \(-\hat d=|c|^2\hat d\), which is impossible for a nonzero operator.

\section{Time-reversal periodic table}

Consider a set of time reversals that all satisfy Eq.~\eqref{eq:physical_TR_conditions} and anticommute pairwise,
\begin{equation}
    \Theta_a\Theta_b=-\Theta_b\Theta_a
    \quad\Longleftrightarrow\quad
    \overline U_bU_a=-\overline U_aU_b
    \qquad (a\ne b).
    \label{eq:anticommuting_TR}
\end{equation}
The two fundamental sign-free sets are the Majorana class \(\{+,-\}\) and the Kramers class \(\{-,-\}\) introduced for Majorana fermions in Ref.~\cite{li2016prl}.  We use the same names for their bosonic extensions because the symmetry algebra and positivity mechanisms are unchanged.  The additional bosonic \(\{+,+,+\}\) set realizes the determinant structure of the split orthogonal group~\cite{wang2015prl}, and we call it the split-orthogonal class.

\begin{center}
    \refstepcounter{table}\label{tab:periodic_table}
    \begin{minipage}{0.96\linewidth}
    \small
    \textbf{TABLE \thetable.}\enspace
    Time-reversal periodic table for bosonic and fermionic Gaussian weights with quadratic evolution that may contain pairing terms.  Every auxiliary-field slice preserves the listed set of pairwise anticommuting antiunitary symmetries.  Each sign records \(\Theta_a^2=\pm I\), so \(\{+,-\}\) contains one symmetry of each square.  The entries list all sets sufficient for \(p_{\rm B}>0\) away from bosonic poles or \(p_{\rm F}\geq0\).  ``All'' covers every assignment of squares.  ``None'' means that these time-reversal symmetries alone do not fix the sign.  The Majorana \(\{+,-\}\) and Kramers \(\{-,-\}\) names follow Ref.~\cite{li2016prl}.  The additional bosonic \(\{+,+,+\}\) entry is the split-orthogonal class, named after the determinant mechanism of Ref.~\cite{wang2015prl}.
    \par\vspace{0.7\baselineskip}
    \centering
    \begin{ruledtabular}
    \begin{tabular}{cll}
        Number of symmetries & \multicolumn{1}{c}{Bosons} & \multicolumn{1}{c}{Fermions} \\
        \hline
        0 or 1 & None & None \\
        2 & Majorana \(\{+,-\}\), Kramers \(\{-,-\}\)
          & Majorana \(\{+,-\}\), Kramers \(\{-,-\}\) \\
        3 & All
          & \(\{+,+,-\},\{+,-,-\},\{-,-,-\}\) \\
        \(\geq4\) & All & All
    \end{tabular}
    \end{ruledtabular}
    \end{minipage}
\end{center}

Three physical mechanisms generate the table.  The Majorana class divides the Nambu operators into two sectors.  Each sector contains annihilation operators together with their Hermitian-conjugate creation operators, and the two sector weights are complex conjugates.  The Kramers class defines a conserved one-body charge, selects one operator from every Nambu-conjugate pair, and leaves Kramers time reversal among the selected operators.  The split-orthogonal class supplies the third bosonic mechanism.  These statements fix the complete Fock-space weights rather than only the determinants under the square roots in Eq.~\eqref{eq:nambu_trace}.

The unlisted two-symmetry class already has a one-mode bosonic counterexample.  Let \(\hat\chi=(\hat b,\hat b^+)^\trans\), \(g=-\ii\sigma_y\), and \(C=\sigma_x\).  The two time reversals
\begin{equation}
    \Theta_1=\mc K,
    \qquad
    \Theta_2=\ii\sigma_z\mc K
    \label{eq:boson_pp_counter_TR}
\end{equation}
obey all three conditions in Eq.~\eqref{eq:physical_TR_conditions}, square to \(+I\), and anticommute.  They both preserve
\begin{equation}
    A=\kappa\sigma_z,
    \qquad \kappa\in\mathbb R,
    \qquad
    \hat h=\kappa\left(\hat n+\frac12\right).
    \label{eq:boson_pp_counter_h}
\end{equation}
For \(\kappa<0\), the ordinary trace converges and equals \(\ee^{\kappa/2}/(1-\ee^\kappa)\).  Its unique analytic continuation to \(\kappa>0\) is
\begin{equation}
    p_{\rm B}=\frac{\ee^{\kappa/2}}{1-\ee^\kappa}
    =-\frac{1}{2\sinh(\kappa/2)}<0.
    \label{eq:boson_pp_counter_weight}
\end{equation}
The occupation-number sum diverges at \(\kappa>0\), while the continued Gaussian weight defined above is finite and negative, with \(\det(I-M)\ne0\).  Thus the bosonic \(\{+,+\}\) set does not fix the sign.  A fermionic counterexample for the \(\{+,+,+\}\) set will be given below.

A single negative-square time reversal is likewise insufficient for general bosonic Nambu evolution.  For two bosonic flavors, let the first Pauli matrix act in Nambu space and the second in flavor space, and take
\begin{equation}
    g=(-\ii\sigma_y)\otimes I_2,
    \quad C=\sigma_x\otimes I_2,
    \quad U=I_2\otimes(\ii\sigma_y),
    \quad \Theta=U\mc K,
    \quad A=x(-\ii\sigma_y)\otimes I_2,
    \label{eq:boson_single_minus_counterexample}
\end{equation}
where \(x\) is real.  The matrix \(\ii\sigma_y\) is real, and direct multiplication gives \(gA+A^\trans g=0\), \(U^\trans gU=g\), \(C\overline UC=U\), \(\overline U U=-I_4\), and \(U^{-1}\overline AU=A\).  The quadratic operator separates into two identical flavors,
\begin{equation}
    \hat h=-\frac{x}{2}\sum_{\alpha=1}^2
      (\hat b_\alpha^2+\hat b_\alpha^{+2}).
\end{equation}
Their continued weights are identical.  For either single flavor, Eq.~\eqref{eq:nambu_trace} gives \(p_1^2=-[4\sin^2(x/2)]^{-1}\).  The two-flavor weight is therefore \(p_{\rm B}=p_1^2\), with no remaining branch choice:
\begin{equation}
    p_{\rm B}=-\frac{1}{4\sin^2(x/2)}<0,
    \qquad x\notin2\pi\mathbb Z.
    \label{eq:boson_single_minus_weight}
\end{equation}
These excluded values are precisely the determinant poles.

\section{Majorana and Kramers classes}

\subsection{Majorana class \texorpdfstring{\(\{+,-\}\)}{\{+,-\}}}

Let
\begin{equation}
    \Theta_+^2=+I,
    \qquad
    \Theta_-^2=-I,
    \qquad
    \Theta_+\Theta_-=-\Theta_-\Theta_+,
    \label{eq:majorana_class}
\end{equation}
and form the linear transformation
\begin{equation}
    P=\Theta_+\Theta_-.
    \label{eq:P_definition}
\end{equation}
The two complex conjugations cancel in this product.  Using Eq.~\eqref{eq:majorana_class},
\begin{equation}
    P^2=-\Theta_+^2\Theta_-^2=I,
    \qquad
    [P,A_\ell]=0,
    \qquad
    P\Theta_+=-\Theta_+P.
    \label{eq:P_properties}
\end{equation}
The first relation divides the one-body operators into \(P=+1\) and \(P=-1\) sectors.  The second keeps these sectors separate throughout the imaginary-time evolution.  Since \(P\) is the product of two physical time reversals, it preserves Hermitian conjugation.  Thus
\begin{equation}
    P\hat d=\pm\hat d
    \quad\Longrightarrow\quad
    P\hat d^+=(P\hat d)^+=\pm\hat d^+.
    \label{eq:P_adjoint_sector}
\end{equation}
The product \(P\) also preserves the canonical matrix, \(P^\trans gP=g\).  To see why this matters, let \(v_+\) and \(v_-\) be the coefficient vectors of one-body operators in the two sectors.  Their canonical pairing vanishes because
\begin{equation}
    v_+^\trans gv_-
    =(Pv_+)^\trans g(Pv_-)
    =-v_+^\trans gv_-
    =0.
    \label{eq:P_sector_canonical_pairing}
\end{equation}
Since \(g\) is invertible, its restriction to either sector is also invertible.  Together with Eq.~\eqref{eq:P_adjoint_sector}, this means that each sector contains a self-contained set of canonical degrees of freedom.  For bosons, the antisymmetry of \(g_{\rm B}\) makes the dimension of each sector even, so the two sectors define separate many-body subsystems.  A fermionic sector can contain an odd number of self-adjoint Majorana components.  In that case, append one decoupled fermionic orbital and assign one of its two Majorana components to each sector.  Both sectors then have even dimension and define separate subsystems.  The added orbital multiplies every trace by two and cannot change its sign.

The last relation in Eq.~\eqref{eq:P_properties} makes \(\Theta_+\) exchange the two sectors:
\begin{equation}
    P\hat d=\pm\hat d
    \quad\Longrightarrow\quad
    P(\Theta_+\hat d)=-\Theta_+(P\hat d)
    =\mp\Theta_+\hat d.
    \label{eq:P_sector_exchange}
\end{equation}
Since \(\Theta_+\) is antiunitary, the complete evolution in one sector is the complex conjugate of that in the other.  If \(p_+\) is the Gaussian trace in the \(P=+1\) sector, the full weight is~\cite{li2015prb,li2016prl}
\begin{equation}
    p=p_+\overline{p_+}=|p_+|^2\geq0.
    \label{eq:majorana_modulus_square}
\end{equation}
Equivalently,
\begin{equation}
    p_{\rm B}=\left|\det\!\left(I-M\big|_{P=+1}\right)\right|^{-1}>0,
    \qquad
    p_{\rm F}=\left|\det\!\left(I+M\big|_{P=+1}\right)\right|\geq0.
    \label{eq:majorana_determinants}
\end{equation}
\subsection{Kramers class \texorpdfstring{\(\{-,-\}\)}{\{-,-\}}}

Now take
\begin{equation}
    \Theta_1^2=\Theta_2^2=-I,
    \qquad
    \Theta_1\Theta_2=-\Theta_2\Theta_1,
    \qquad
    Q=\Theta_1\Theta_2.
    \label{eq:kramers_class}
\end{equation}
The product \(Q\) is linear and obeys
\begin{equation}
    Q^2=-I,
    \qquad
    [Q,A_\ell]=0.
    \label{eq:Q_properties}
\end{equation}
Its eigenvalues are \(+\ii\) and \(-\ii\).  Hermitian conjugation exchanges these two eigenspaces because \(Q\) is a physical unitary transformation:
\begin{equation}
    Q\hat d=\ii\hat d
    \quad\Longrightarrow\quad
    Q\hat d^+=(Q\hat d)^+=-\ii\hat d^+.
    \label{eq:Q_adjoint_pair}
\end{equation}
The \(Q=+\ii\) sector therefore contains exactly \(N\) independent one-body operators, while the \(Q=-\ii\) sector contains their Hermitian conjugates.  The matrix \(\ii Q\) generates a conserved quadratic charge because
\begin{equation}
    \left[\frac12\hat\chi^\trans g(\ii Q)\hat\chi,
      \hat h_\ell\right]
    =\frac12\hat\chi^\trans g[\ii Q,A_\ell]\hat\chi=0.
    \label{eq:charge_conservation}
\end{equation}
Thus \(Q\) supplies the particle-number-like structure that a general Nambu problem lacks.

Choose a column \(\hat{\bb d}=(\hat d_1,\ldots,\hat d_N)^\trans\) of one-body operators spanning the \(Q=+\ii\) sector.  We define the reduced slice generators and their ordered product by
\begin{equation}
    \ee^{-\hat h_\ell}\hat{\bb d}\,\ee^{\hat h_\ell}
    =\ee^{A_{\ell,+\ii}}\hat{\bb d},
    \qquad
    M_{+\ii}=\ee^{A_{1,+\ii}}\cdots\ee^{A_{L,+\ii}}.
    \label{eq:reduced_matrices}
\end{equation}
Normalize the independent fermionic operators so that \(\{\hat d_\mu,\hat d_\nu^+\}=D_{\mu\nu}\).  For bosons, choose a basis in which \([\hat d_\mu,\hat d_\nu^+]=D_{\mu\nu}\) is diagonal.  In the two cases,
\begin{equation}
    D=I_N\quad\text{for fermions},
    \qquad
    D=\diag(I_{N-N_-},-I_{N_-})\quad\text{for bosons}.
    \label{eq:reduced_commutator}
\end{equation}
For bosons, \(N_-\) counts the selected operators whose Hermitian conjugates are the physical annihilation operators.  In a basis formed by the \(N\) selected operators followed by their Hermitian conjugates, every generator and the full propagation have the forms
\begin{equation}
    A_\ell\sim
    \begin{pmatrix}
        A_{\ell,+\ii}&0\\
        0&-D A_{\ell,+\ii}^\trans D
    \end{pmatrix},
    \qquad
    M\sim
    \begin{pmatrix}
        M_{+\ii}&0\\
        0&D(M_{+\ii}^\trans)^{-1}D
    \end{pmatrix}.
    \label{eq:Q_block_propagation}
\end{equation}
Consequently,
\begin{align}
    (-1)^N\det(I_{2N}-M)
    &=\det(M_{+\ii})^{-1}\det(I_N-M_{+\ii})^2,
    \label{eq:boson_reduced_square}\\
    \det(I_{2N}+M)
    &=\det(M_{+\ii})^{-1}\det(I_N+M_{+\ii})^2,
    \label{eq:fermion_reduced_square}\\
    \det(M_{+\ii})
    &=\exp\!\left[\sum_\ell\Tr A_{\ell,+\ii}\right].
    \label{eq:reduced_det_exp}
\end{align}
These identities determine the magnitude of the weight.  Its sign follows from how the \(Q=+\ii\) operators sit in Fock space.  For fermions, the selected operators are canonical annihilation operators \(\hat d_\mu\).  Inserting the first block matrix in Eq.~\eqref{eq:Q_block_propagation} into Eq.~\eqref{eq:configuration_weight} gives
\begin{equation}
    \hat h_\ell
    =\sum_{\mu\nu}\hat d_\mu^+
      (A_{\ell,+\ii})_{\mu\nu}\hat d_\nu
     -\frac12\Tr A_{\ell,+\ii}.
    \label{eq:fermion_reduced_operator}
\end{equation}
For one eigenvalue \(x\), the trace contributes \(\ee^{-x/2}(1+\ee^x)\), as in Eq.~\eqref{eq:fermion_one_mode}.  The usual fermionic Gaussian trace over all \(N\) independent operators therefore gives
\begin{equation}
    p_{\rm F}=
    \exp\!\left[-\frac12\sum_\ell\Tr A_{\ell,+\ii}\right]
    \det(I_N+M_{+\ii}).
    \label{eq:fermion_reduced_weight}
\end{equation}
For bosons, the same substitution gives
\begin{equation}
    \hat h_\ell
    =\sum_{\mu\nu}\hat d_\mu^+
      (D A_{\ell,+\ii})_{\mu\nu}\hat d_\nu
     +\frac12\Tr A_{\ell,+\ii}.
    \label{eq:boson_reduced_operator}
\end{equation}
The sign is visible at a stable diagonal reference point, where the selected directions do not mix.  If the selected operator is an annihilation operator and its propagation eigenvalue is \(\ee^x\), the trace contributes \(\ee^{x/2}(1-\ee^x)^{-1}\).  Selecting the creation operator reverses \(x\) and contributes
\begin{equation}
    \frac{\ee^{-x/2}}{1-\ee^{-x}}
    =-\frac{\ee^{x/2}}{1-\ee^x}.
    \label{eq:boson_reversed_mode}
\end{equation}
Thus the reference trace carries one minus sign for every negative entry of \(D\).  A bosonic coherent-state Gaussian integral evaluates Eq.~\eqref{eq:boson_reduced_operator} when the positive and negative directions mix and when different slices do not commute~\cite{munn1978ensemble}.  The integral and the expression below agree on a nonempty open set of trace-convergent generators.  They are the same meromorphic function of the generator matrix elements and therefore agree at every point where \(\det(I_N-M_{+\ii})\ne0\).  The exact continued weight is
\begin{equation}
    p_{\rm B}=(-1)^{N_-}
    \exp\!\left[\frac12\sum_\ell\Tr A_{\ell,+\ii}\right]
    \det(I_N-M_{+\ii})^{-1}.
    \label{eq:boson_reduced_weight}
\end{equation}

Although \(Q\Theta_1=-\Theta_1Q\), the anti-linearity of \(\Theta_1\) keeps the \(Q=+\ii\) sector invariant:
\begin{equation}
    Q\hat d=\ii\hat d
    \quad\Longrightarrow\quad
    Q(\Theta_1\hat d)=-\Theta_1(Q\hat d)
    =-\Theta_1(\ii\hat d)=\ii\Theta_1\hat d.
    \label{eq:Kramers_in_reduced_space}
\end{equation}
Hence \(\Theta_1^2=-I\) acts as ordinary Kramers time reversal on the independent operators.  If an eigenoperator of the complete one-body evolution obeys
\begin{equation}
    \ee^{-\hat h_L}\cdots\ee^{-\hat h_1}
    \hat d\,
    \ee^{\hat h_1}\cdots\ee^{\hat h_L}
    =\lambda\hat d,
    \quad\text{then}\quad
    \ee^{-\hat h_L}\cdots\ee^{-\hat h_1}
    (\Theta_1\hat d)
    \ee^{\hat h_1}\cdots\ee^{\hat h_L}
    =\overline\lambda\,\Theta_1\hat d.
    \label{eq:Kramers_eigenvalues}
\end{equation}
Nonreal eigenvalues appear in complex-conjugate pairs.  For a real eigenvalue, \(\hat d\) and \(\Theta_1\hat d\) are linearly independent because \(\Theta_1^2=-I\).  The same pairing holds for every vector in a Jordan chain, so each real eigenvalue has even algebraic multiplicity.  It follows that
\begin{equation}
    \det(I_N\mp M_{+\ii})
    =\prod_{\lambda\notin\mathbb R}|1\mp\lambda|^2
     \prod_{\lambda\in\mathbb R}(1\mp\lambda)^2
    \geq0,
    \label{eq:Kramers_determinant}
\end{equation}
where each conjugate or Kramers pair is included once.  The same symmetry gives \(\Tr A_{\ell,+\ii}\in\mathbb R\).

For bosons, the negative entries in Eq.~\eqref{eq:reduced_commutator} also form Kramers pairs.  For any selected operator \(\hat d\), time reversal gives
\begin{equation}
    [\Theta_1\hat d,(\Theta_1\hat d)^+]
    =\overline{[\hat d,\hat d^+]}
    =[\hat d,\hat d^+],
    \qquad
    [\hat d,(\Theta_1\hat d)^+]=0.
    \label{eq:commutator_Kramers_pair}
\end{equation}
The first commutator is real because it is Hermitian.  To prove the second relation, denote \(z=[\hat d,(\Theta_1\hat d)^+]\).  Hermitian conjugation gives \(\overline z=[\Theta_1\hat d,\hat d^+]\).  Applying time reversal to the definition of \(z\) instead gives
\begin{equation}
    \overline z
    =[\Theta_1\hat d,(\Theta_1^2\hat d)^+]
    =-[\Theta_1\hat d,\hat d^+]
    =-\overline z,
\end{equation}
so \(z=0\).  Every selected creation-like operator therefore has an independent, orthogonal partner with the same negative commutator, and \(N_-\) is even.  Equations~\eqref{eq:fermion_reduced_weight}--\eqref{eq:Kramers_determinant} now give
\begin{equation}
    \boxed{p_{\rm F}\geq0,\qquad p_{\rm B}>0\ \text{away from bosonic poles}.}
    \label{eq:kramers_class_positive}
\end{equation}
This is the Nambu-space form of the Kramers positivity mechanism~\cite{li2016prl}.

For a particle-number-conserving problem, the annihilation-operator space is already an independent \(N\)-dimensional space.  A single negative-square time reversal then produces Eq.~\eqref{eq:Kramers_eigenvalues} directly, which is the familiar Kramers criterion for determinant quantum Monte Carlo~\cite{wu2005prb}.  The second time reversal in Eq.~\eqref{eq:kramers_class} supplies this independent space when pairing mixes annihilation and creation operators.

\section{Three time reversals and the \texorpdfstring{\(\{+,+,+\}\)}{\{+,+,+\}} distinction}

Three sign patterns already contain one of the elementary pairs.  The classes \(\{+,+,-\}\) and \(\{+,-,-\}\) contain a \(\{+,-\}\) pair, while \(\{-,-,-\}\) contains a \(\{-,-\}\) pair.  They are sign-problem-free for both bosons and fermions.  The remaining class \(\{+,+,+\}\) distinguishes the two statistics.

\subsection{Bosonic split-orthogonal class \texorpdfstring{\(\{+,+,+\}\)}{\{+,+,+\}}}

For bosons, a positive-square time reversal defines a real set of canonical Nambu operators.  To construct it, start from any operator \(\hat d\) and form
\begin{equation}
    \hat d+\widehat\Theta_1\hat d\widehat\Theta_1^{-1},
    \qquad
    \ii(\hat d-\widehat\Theta_1\hat d\widehat\Theta_1^{-1}).
    \label{eq:Theta1_real_operators}
\end{equation}
Both combinations are fixed by \(\Theta_1\).  They span the full Nambu space over the real numbers.  Since \(\Theta_1\) preserves bosonic commutators and Hermitian conjugation, these operators can be normalized in canonical conjugate pairs while staying inside this real span.  In the resulting basis, the three symmetries of the \(\{+,+,+\}\) class have the form
\begin{equation}
    \Theta_1=\mc K,
    \qquad
    \Theta_2=\ii R_2\mc K,
    \qquad
    \Theta_3=\ii R_3\mc K,
    \label{eq:three_plus_basis}
\end{equation}
with real matrices \(R_2\) and \(R_3\).  The first equality follows from the construction above.  Anticommutation with \(\Theta_1=\mc K\) makes the other two one-body matrices purely imaginary, which accounts for the factors of \(\ii\).  Substituting Eq.~\eqref{eq:three_plus_basis} into Eq.~\eqref{eq:physical_TR_conditions} gives
\begin{gather}
    \overline{A_\ell}=A_\ell,
    \qquad
    [R_2,A_\ell]=[R_3,A_\ell]=0,
    \qquad
    R_2R_3=-R_3R_2,
    \label{eq:three_plus_A}\\
    R_j^2=I,
    \qquad
    R_j^\trans gR_j=-g,
    \qquad
    CR_jC=-R_j,
    \qquad j=2,3.
    \label{eq:three_plus_physical}
\end{gather}
The first two time reversals define
\begin{equation}
    Q=\Theta_1\Theta_2=-\ii R_2,
    \qquad Q^2=-I.
    \label{eq:three_plus_Q}
\end{equation}
We again work in its \(Q=+\ii\) sector.  All \(A_{\ell,+\ii}\) and \(M_{+\ii}\) are real in this basis.  The third time reversal supplies the real matrix
\begin{equation}
    \eta=(gR_3)\big|_{Q=+\ii}.
    \label{eq:eta_definition}
\end{equation}
It is invertible on this sector.  Indeed, \(Q^\trans(gR_3)Q=-gR_3\).  If \(Qv=+\ii v\) and \(Qw=-\ii w\), this identity gives \(v^\trans gR_3w=-v^\trans gR_3w=0\).  Thus \(gR_3\) has no matrix elements between the two sectors and remains invertible inside each one.

From \(R_3^\trans gR_3=-g\) and \(R_3^2=I\),
\begin{equation}
    R_3^\trans g=-gR_3.
    \label{eq:R3g_relation}
\end{equation}
Together with \(g_{\rm B}^\trans=-g_{\rm B}\), this makes \(\eta\) symmetric:
\begin{equation}
    \eta^\trans=\eta.
    \label{eq:eta_symmetric}
\end{equation}
Every reduced generator preserves this matrix because
\begin{align}
    A_{\ell,+\ii}^\trans\eta+\eta A_{\ell,+\ii}
    &=\left.g(R_3A_\ell-A_\ell R_3)\right|_{Q=+\ii}
    \notag\\
    &=\left.g[R_3,A_\ell]\right|_{Q=+\ii}=0.
    \label{eq:eta_generator_relation}
\end{align}

In this canonical Nambu basis, the physical commutator matrix is
\begin{equation}
    [\hat\chi_i,\hat\chi_j^+]=(-gC)_{ij}.
    \label{eq:boson_commutator_matrix}
\end{equation}
Its restriction to the \(Q=+\ii\) operators becomes the matrix \(D\) in Eq.~\eqref{eq:reduced_commutator} after the selected operators are normalized.
The matrix
\begin{equation}
    S=(CR_3)\big|_{Q=+\ii}
    \label{eq:S_definition}
\end{equation}
stays within the \(Q=+\ii\) sector because \(R_2\) commutes with \(CR_3\).  Using \(CgC=-g\), \(CR_3C=-R_3\), and Eq.~\eqref{eq:R3g_relation}, one finds
\begin{align}
    S^\trans\eta S
    &=R_3^\trans(CgR_3C)R_3
      =R_3^\trans g=-gR_3=-\eta,
    \notag\\
    S^\trans\left.(-gC)\right|_{Q=+\ii}S
    &=-R_3^\trans CgR_3
      =-\left.(-gC)\right|_{Q=+\ii}.
    \label{eq:S_flips_forms}
\end{align}
An invertible change of basis cannot alter how many positive and negative directions a real symmetric matrix has.  The first equality says that \(S\) turns \(\eta\) into \(-\eta\), so those two numbers must be equal.  The second equality gives the same result for the commutator in Eq.~\eqref{eq:boson_commutator_matrix}.  Hence \(N\) is even, both matrices have \(N/2\) positive and \(N/2\) negative directions, and the integer in Eq.~\eqref{eq:boson_reduced_weight} is fixed to
\begin{equation}
    N_-=\frac N2.
    \label{eq:Nminus_half}
\end{equation}

After a real change of basis, take
\begin{equation}
    \eta=\diag(I_{N/2},-I_{N/2}).
    \label{eq:eta_diagonal}
\end{equation}
The two diagonal blocks of \(\eta\) label equal sets of positive and negative channels.  Equation~\eqref{eq:eta_generator_relation} implies
\begin{equation}
    M_{+\ii}^\trans\eta M_{+\ii}=\eta,
    \qquad
    \Tr A_{\ell,+\ii}=0.
    \label{eq:eta_propagation}
\end{equation}
Real matrices satisfying the first relation form the split orthogonal group \(O(N/2,N/2)\).  Each factor \(\ee^{A_{\ell,+\ii}}\) can be reached continuously from the identity along \(\ee^{tA_{\ell,+\ii}}\) while preserving \(\eta\).  The ordered product therefore lies in the same connected part of the group as the identity.  The split-orthogonal determinant theorem~\cite{wang2015prl} gives
\begin{equation}
    \det(I_N+M_{+\ii})\geq0.
    \label{eq:plus_determinant_positive}
\end{equation}
The eigenvalue structure explains the sign.  Preservation of \(\eta\) produces reciprocal eigenvalues, while reality produces complex-conjugate eigenvalues.  Away from the unit circle, a nonreal quartet \(\lambda,\overline\lambda,\lambda^{-1},\overline\lambda^{-1}\) contributes \(|1+\lambda|^2|1+\lambda^{-1}|^2\).  On the unit circle, the quartet reduces to the conjugate pair \(\lambda,\overline\lambda\), which contributes \(|1+\lambda|^2\).  A positive real reciprocal pair contributes \((1+\lambda)(1+\lambda^{-1})>0\).  In the component connected to the identity, negative real reciprocal pairs occur in an even number, so their negative contributions also multiply to a positive number.

The bosonic trace contains \(I_N-M_{+\ii}\).  Its sign follows without another spectral classification.  Whenever \(I_N+M_{+\ii}\) is invertible, Eq.~\eqref{eq:eta_propagation} gives
\begin{equation}
    \left[\eta(I_N-M_{+\ii})(I_N+M_{+\ii})^{-1}\right]^\trans
    =-\eta(I_N-M_{+\ii})(I_N+M_{+\ii})^{-1}.
    \label{eq:cayley_antisymmetric}
\end{equation}
The matrix in brackets is real and antisymmetric, so its determinant is the square of its Pfaffian.  Since \(\det\eta=(-1)^{N/2}\),
\begin{equation}
    (-1)^{N/2}
    \frac{\det(I_N-M_{+\ii})}{\det(I_N+M_{+\ii})}
    =\operatorname{Pf}\!\left[
      \eta(I_N-M_{+\ii})(I_N+M_{+\ii})^{-1}
    \right]^2\geq0.
    \label{eq:minus_determinant_positive}
\end{equation}
The result extends by continuity when \(I_N+M_{+\ii}\) is singular.  Combining Eqs.~\eqref{eq:plus_determinant_positive} and~\eqref{eq:minus_determinant_positive} yields
\begin{equation}
    (-1)^{N/2}\det(I_N-M_{+\ii})\geq0.
    \label{eq:boson_three_plus_det_sign}
\end{equation}
Finally, Eqs.~\eqref{eq:boson_reduced_weight}, \eqref{eq:Nminus_half}, and~\eqref{eq:eta_propagation} give
\begin{equation}
    \boxed{
    p_{\rm B}=(-1)^{N/2}\det(I_N-M_{+\ii})^{-1}>0
    }
    \label{eq:boson_three_plus_weight}
\end{equation}
away from bosonic poles.  The bosonic split-orthogonal class \(\{+,+,+\}\) therefore protects the full Gaussian weight.

\subsection{Fermionic \texorpdfstring{\(\{+,+,+\}\)}{\{+,+,+\}} class}

The fermionic \(\{+,+,+\}\) class does not protect a general Nambu evolution.  A two-mode counterexample makes the difference explicit.  Let the first Pauli matrix act in Nambu space and the second in flavor space, and choose
\begin{gather}
    g=C=\sigma_x\otimes I_2,
    \qquad
    U_1=I_4,
    \quad
    U_2=\ii\sigma_z\otimes\sigma_x,
    \quad
    U_3=\ii\sigma_z\otimes\sigma_z,
    \label{eq:fermion_three_plus_U}\\
    A_1=\pi\ii\sigma_x\otimes\sigma_y,
    \qquad
    A_2=x\sigma_z\otimes I_2,
    \qquad x\in\mathbb R.
    \label{eq:fermion_three_plus_A}
\end{gather}
The Pauli algebra gives, for \(a\ne b\) and \(\ell=1,2\),
\begin{equation}
\begin{gathered}
    gA_\ell+A_\ell^\trans g=0,
    \qquad \overline U_aU_a=I_4,
    \qquad \overline U_bU_a=-\overline U_aU_b,
    \qquad U_a^\trans gU_a=g,\\
    C\overline U_aC=U_a,
    \qquad U_a^{-1}\overline{A_\ell}U_a=A_\ell.
\end{gathered}
    \label{eq:fermion_counter_checks}
\end{equation}
Thus \(\Theta_a=U_a\mc K\) are three physical, pairwise anticommuting time reversals of both slices, and every one squares to \(+I\).  In Fock space,
\begin{equation}
    \hat h_1=\pi(\hat c_1\hat c_2+\hat c_1^+\hat c_2^+),
    \qquad
    \hat h_2=x(\hat n_1+\hat n_2-1).
    \label{eq:fermion_counter_operators}
\end{equation}
The first exponential is \(+I\) in the odd-parity sector and \(-I\) in the even-parity sector.  The two odd states contribute \(2\), while the empty and doubly occupied states contribute \(-\ee^{-x}\) and \(-\ee^x\).  Therefore
\begin{equation}
    p_{\rm F}=\Tr(\ee^{\hat h_1}\ee^{\hat h_2})
    =2-\ee^x-\ee^{-x}<0
    \qquad(x\in\mathbb R\setminus\{0\}).
    \label{eq:fermion_three_plus_negative}
\end{equation}
The fermionic \(\{+,+,+\}\) class consequently has no symmetry-protected sign.  The product \(\Theta_1\Theta_2\Theta_3\) is itself a physical time reversal, preserves both slices, and squares to \(-I\).  This same counterexample also shows why a single negative-square time reversal is insufficient for a general paired fermion problem.

\section{Four or more time reversals}

Four mutually anticommuting time reversals always generate an elementary sign-free pair.  If both square signs occur, one can select a \(\{+,-\}\) pair immediately.  If all four squares have the same sign, form the product of three,
\begin{equation}
    \widetilde\Theta=\Theta_2\Theta_3\Theta_4.
    \label{eq:generated_TR}
\end{equation}
This product is antiunitary and preserves every quadratic slice.  Pairwise anticommutation gives
\begin{equation}
    \widetilde\Theta^2
    =-\Theta_2^2\Theta_3^2\Theta_4^2,
    \qquad
    \Theta_1\widetilde\Theta
    =-\widetilde\Theta\Theta_1.
    \label{eq:generated_TR_properties}
\end{equation}
The square of \(\widetilde\Theta\) is opposite to that of \(\Theta_1\), so the generated pair is of type \(\{+,-\}\).  Any larger anticommuting set contains four members.  Thus four or more time reversals guarantee nonnegative bosonic and fermionic weights.

\section{Bosonic models with anomalous propagation}

The time-reversal criteria act on every fixed auxiliary-field slice and therefore remain effective when the quadratic evolution contains pairing terms.  The following Hermitian interacting models illustrate the \(\{+,-\}\) and \(\{+,+,+\}\) classes.

\subsection{Conjugate Bogoliubov copies: Majorana class \texorpdfstring{\(\{+,-\}\)}{\{+,-\}}}

Place two bosonic flavors \(\hat b_i\) and \(\hat c_i\) on each site and consider
\begin{align}
    \hat H_0={}&
    \sum_{ij}\left(
      \hat b_i^+K_{ij}\hat b_j
      +\hat c_i^+\overline{K}_{ij}\hat c_j
    \right)
    \notag\\
    &+\frac12\sum_{ij}\left(
      \Delta_{ij}\hat b_i^+\hat b_j^+
      +\overline\Delta_{ij}\hat c_i^+\hat c_j^+
      +\overline\Delta_{ij}\hat b_j\hat b_i
      +\Delta_{ij}\hat c_j\hat c_i
    \right),
    \label{eq:conjugate_BdG_H0}\\
    \hat H_I={}&U\sum_i(\hat n_{b,i}-\hat n_{c,i})^2,
    \qquad U>0,
    \label{eq:conjugate_BdG_HI}
\end{align}
where \(K=K^\dagger\) and \(\Delta=\Delta^\trans\).  A sufficient stability condition for the quadratic part is
\begin{equation}
    \begin{pmatrix}K&\Delta\\ \Delta^\dagger&K^\trans\end{pmatrix}>0.
    \label{eq:conjugate_BdG_stability}
\end{equation}
On a finite lattice this condition makes \(\ee^{-\Delta\tau\hat H_0}\) damp every large-occupation state.  The auxiliary-field factors below are unitary.  Their ordered products consequently have finite many-body traces, so the model weights do not encounter bosonic poles.

On the flavor doublet, define
\begin{equation}
    \widehat\Theta_+
    \begin{pmatrix}\hat b_i\\\hat c_i\end{pmatrix}
    \widehat\Theta_+^{-1}
    =\sigma_x\begin{pmatrix}\hat b_i\\\hat c_i\end{pmatrix},
    \qquad
    \widehat\Theta_-
    \begin{pmatrix}\hat b_i\\\hat c_i\end{pmatrix}
    \widehat\Theta_-^{-1}
    =\ii\sigma_y\begin{pmatrix}\hat b_i\\\hat c_i\end{pmatrix}.
    \label{eq:conjugate_BdG_TR}
\end{equation}
Both transformations conjugate numerical coefficients.  Their one-body operator maps obey
\begin{equation}
    \Theta_+^2=+I,
    \qquad
    \Theta_-^2=-I,
    \qquad
    \Theta_+\Theta_-=-\Theta_-\Theta_+.
    \label{eq:conjugate_BdG_TR_algebra}
\end{equation}
The Fock-space transformations preserve \(\hat H_0\) and send
\begin{equation}
    \hat n_{b,i}-\hat n_{c,i}
    \longmapsto
    -(\hat n_{b,i}-\hat n_{c,i}).
    \label{eq:relative_density_odd}
\end{equation}
The continuous Hubbard--Stratonovich identity
\begin{equation}
    \ee^{-\Delta\tau U(\hat n_{b,i}-\hat n_{c,i})^2}
    =\frac{1}{\sqrt{2\pi}}
    \int_{-\infty}^{\infty}d\phi_i\,\ee^{-\phi_i^2/2}
    \ee^{\ii\sqrt{2\Delta\tau U}\,\phi_i
      (\hat n_{b,i}-\hat n_{c,i})}
    \label{eq:relative_density_HS}
\end{equation}
therefore preserves both time reversals for every real value of \(\phi_i\).  Complex conjugation changes the sign of \(\ii\), while Eq.~\eqref{eq:relative_density_odd} supplies the second sign change.

At fixed auxiliary fields, the \(b\) and \(c\) evolutions remain separate and are exact complex conjugates.  Their many-body traces obey
\begin{equation}
    p(\phi)=p_b(\phi)\overline{p_b(\phi)}
    =|p_b(\phi)|^2\geq0.
    \label{eq:conjugate_BdG_weight}
\end{equation}
The matrix \(K\) may describe frustrated hopping or flux, while \(\Delta\) may describe two chiral pairing patterns such as \(d+\ii d\) and \(d-\ii d\).  Equation~\eqref{eq:conjugate_BdG_weight} makes interaction-driven transitions between paired, density-imbalanced, and phase-coherent bosonic states accessible to auxiliary-field simulation on frustrated lattices.

\subsection{Frustrated triangular-lattice two-flavor \texorpdfstring{\(p\)}{p}-wave model in the split-orthogonal class \texorpdfstring{\(\{+,+,+\}\)}{\{+,+,+\}}}

Frustrated hopping on the triangular lattice produces two low-energy valleys at \(\pm\mathbf K\), allowing orbital motion to become locked to an internal flavor.  For two-component bosons, closely related multivalley settings have been predicted to support counterpropagating flavors with no net mass current and staggered spin loop currents.  When the valley bands carry nonzero Berry curvature, the same order also produces a spontaneous spin Hall response~\cite{liNatuParamekanti2014chiral}.  Density-matrix renormalization-group calculations on a triangular ladder found that the chiral-spin structure survives strong one-dimensional quantum fluctuations~\cite{wuLiDasSarma2017chiral}.  Odd-parity flavor-singlet pairing adds a second source of strong quantum fluctuations and connects this double-valley problem to exactly solvable \(p\)-wave boson models with fragmented pair condensation and an unusual third-order quantum transition~\cite{lermaHernandezDukelsky2019pwave}.  The Hamiltonian below brings double-valley counterflow and \(p\)-wave pair dynamics together in a local two-dimensional model whose auxiliary-field weight is strictly positive.

Let the primitive vectors be \(\mathbf a_1=(1,0)\) and \(\mathbf a_2=(1/2,\sqrt{3}/2)\).  We sum over the three displacements
\begin{equation}
    \boldsymbol\delta\in
    \{\mathbf a_1,\mathbf a_2,\mathbf a_2-\mathbf a_1\},
    \label{eq:three_plus_bond_directions}
\end{equation}
so that every nearest-neighbor bond appears once.  The chosen directions are a physical part of the pairing texture.  On such a bond, define the flavor-singlet pair
\begin{equation}
    \hat P_{\mathbf r,\boldsymbol\delta}^+
    =\hat b_{\mathbf r}^+\hat c_{\mathbf r+\boldsymbol\delta}^+
     -\hat b_{\mathbf r+\boldsymbol\delta}^+
      \hat c_{\mathbf r}^+,
    \qquad
    \hat P_{\mathbf r,\boldsymbol\delta}
    =\hat c_{\mathbf r+\boldsymbol\delta}\hat b_{\mathbf r}
     -\hat c_{\mathbf r}\hat b_{\mathbf r+\boldsymbol\delta}.
    \label{eq:odd_pair}
\end{equation}
Its spatial and flavor wave functions are both antisymmetric.  Their product has the exchange symmetry required for two bosons.  The bond field is odd under \(\boldsymbol\delta\mapsto-\boldsymbol\delta\) and has a linear long-wavelength harmonic, which is the lattice analogue of \(p\)-wave pairing.  This directional texture lowers the threefold rotation symmetry and remains finite at the two valleys.

We consider
\begin{align}
    \hat H={}&\hat H_0+\hat H_I,
    \label{eq:three_plus_model}\\
    \hat H_0={}&
    -\mu\sum_{\mathbf r}
    (\hat n_{b,\mathbf r}+\hat n_{c,\mathbf r})
    +t\sum_{\mathbf r,\boldsymbol\delta}
    \left(
      \hat b_{\mathbf r}^+\hat b_{\mathbf r+\boldsymbol\delta}
      +\hat c_{\mathbf r}^+\hat c_{\mathbf r+\boldsymbol\delta}
      +\mathrm{H.c.}
    \right)
    \notag\\
    &-\Delta_{\rm p}\sum_{\mathbf r,\boldsymbol\delta}
    \left(
      \hat P_{\mathbf r,\boldsymbol\delta}^+
      +\hat P_{\mathbf r,\boldsymbol\delta}
    \right),
    \label{eq:three_plus_model_H0}\\
    \hat H_I={}&
    g_1\sum_{\mathbf r,\boldsymbol\delta}
    \left[
      -\ii\left(
        \hat b_{\mathbf r}^+\hat b_{\mathbf r+\boldsymbol\delta}
        -\hat b_{\mathbf r+\boldsymbol\delta}^+\hat b_{\mathbf r}
        +\hat c_{\mathbf r}^+\hat c_{\mathbf r+\boldsymbol\delta}
        -\hat c_{\mathbf r+\boldsymbol\delta}^+\hat c_{\mathbf r}
      \right)
    \right]^2
    \notag\\
    &+g_2\sum_{\mathbf r,\boldsymbol\delta}
    \left[
      -\ii\left(
        \hat P_{\mathbf r,\boldsymbol\delta}^+
        -\hat P_{\mathbf r,\boldsymbol\delta}
      \right)
    \right]^2,
    \qquad g_1,g_2>0.
    \label{eq:three_plus_model_HI}
\end{align}
Here \(t>0\), while \(\Delta_{\rm p}\) and \(\mu\) are real.  The plus sign of the hopping frustrates every elementary triangle.  In the absence of pairing, the dispersion is
\begin{equation}
    \xi_{\mathbf k}=-\mu+2t\left[
      \cos(\mathbf k\cdot\mathbf a_1)
      +\cos(\mathbf k\cdot\mathbf a_2)
      +\cos\bigl(\mathbf k\cdot(\mathbf a_2-\mathbf a_1)\bigr)
    \right],
    \label{eq:three_plus_dispersion}
\end{equation}
with minima at \(\pm\mathbf K=\pm(4\pi/3,0)\).  The pair term has the odd form factor
\begin{equation}
    2\ii\Delta_{\rm p}\left[
      \sin(\mathbf k\cdot\mathbf a_1)
      +\sin(\mathbf k\cdot\mathbf a_2)
      +\sin\bigl(\mathbf k\cdot(\mathbf a_2-\mathbf a_1)\bigr)
    \right].
    \label{eq:three_plus_pair_form_factor}
\end{equation}
It changes sign between the valleys and pairs a \(b\) boson near one valley with a \(c\) boson near the other.  We choose \(\mu\) sufficiently negative that the free Bogoliubov Hamiltonian is stable.  A simple sufficient condition is that \(\xi_{\mathbf k}\) exceed the magnitude of Eq.~\eqref{eq:three_plus_pair_form_factor} for every \(\mathbf k\).

The model has an exact global flavor \(SU(2)\) symmetry,
\begin{equation}
    \begin{pmatrix}\hat b_{\mathbf r}\\ \hat c_{\mathbf r}\end{pmatrix}
    \longmapsto
    V\begin{pmatrix}\hat b_{\mathbf r}\\ \hat c_{\mathbf r}\end{pmatrix},
    \qquad V\in SU(2).
    \label{eq:three_plus_SU2}
\end{equation}
The hopping and the total current in Eq.~\eqref{eq:three_plus_model_HI} are flavor scalars, while Eq.~\eqref{eq:odd_pair} is an \(SU(2)\) singlet.  In particular, the relative \(U(1)\) transformation
\begin{equation}
    \hat b_{\mathbf r}\longmapsto \ee^{\ii\theta}\hat b_{\mathbf r},
    \qquad
    \hat c_{\mathbf r}\longmapsto \ee^{-\ii\theta}\hat c_{\mathbf r}
    \label{eq:three_plus_relative_U1}
\end{equation}
is an exact subgroup, and
\begin{equation}
    \left[\hat H,\sum_{\mathbf r}
    (\hat n_{b,\mathbf r}-\hat n_{c,\mathbf r})\right]=0.
    \label{eq:three_plus_relative_number}
\end{equation}
The common boson parity
\begin{equation}
    (-1)^{\sum_{\mathbf r}(\hat n_{b,\mathbf r}+\hat n_{c,\mathbf r})}
    =\ee^{\ii\pi\sum_{\mathbf r}
    (\hat n_{b,\mathbf r}-\hat n_{c,\mathbf r})}
    \label{eq:three_plus_common_parity}
\end{equation}
is the central \(-I\) element of this same \(SU(2)\).  It therefore does not supply a separate Ising degree of freedom.

The two interactions have direct physical meanings.  For either flavor on one bond,
\begin{align}
    &\left[-\ii\left(
      \hat b_{\mathbf r}^+\hat b_{\mathbf r+\boldsymbol\delta}
      -\hat b_{\mathbf r+\boldsymbol\delta}^+\hat b_{\mathbf r}
    \right)\right]^2
    \notag\\
    &\quad=
    \hat n_{b,\mathbf r}+\hat n_{b,\mathbf r+\boldsymbol\delta}
    +2\hat n_{b,\mathbf r}\hat n_{b,\mathbf r+\boldsymbol\delta}
    -\left[
      (\hat b_{\mathbf r}^+)^2\hat b_{\mathbf r+\boldsymbol\delta}^2
      +(\hat b_{\mathbf r+\boldsymbol\delta}^+)^2\hat b_{\mathbf r}^2
    \right].
    \label{eq:three_plus_current_square_expansion}
\end{align}
Thus \(g_1\) combines a nearest-neighbor density cost with coherent pair hopping.  The cross term couples the currents of the two flavors.  Positive \(g_1\) penalizes net mass flow and makes counterpropagating flavors more favorable than two flavors moving through the same valley.  Atom-pair hopping has been studied directly in a constrained triangular-lattice Bose--Hubbard model, where large-scale quantum Monte Carlo simulations found pair-superfluid and pair-supersolid phases~\cite{zhangYinWang2013pair}.  The additional density and cross-flavor current terms make Eq.~\eqref{eq:three_plus_model_HI} a distinct, flavor-symmetric realization of correlated pair motion.

The second square obeys
\begin{align}
    &\left[-\ii\left(
      \hat P_{\mathbf r,\boldsymbol\delta}^+
      -\hat P_{\mathbf r,\boldsymbol\delta}
    \right)\right]^2
    \notag\\
    &\quad=
    \hat P_{\mathbf r,\boldsymbol\delta}^+
      \hat P_{\mathbf r,\boldsymbol\delta}
    +\hat P_{\mathbf r,\boldsymbol\delta}
      \hat P_{\mathbf r,\boldsymbol\delta}^+
    -(\hat P_{\mathbf r,\boldsymbol\delta}^+)^2
    -\hat P_{\mathbf r,\boldsymbol\delta}^2.
    \label{eq:three_plus_pair_square_expansion}
\end{align}
It contains number-conserving scattering of odd-parity pairs together with processes that create or remove two bond pairs.  The real linear pair field in Eq.~\eqref{eq:three_plus_model_H0} selects one pair quadrature, while positive \(g_2\) suppresses fluctuations in the orthogonal quadrature.  The exactly solvable model of Ref.~\cite{lermaHernandezDukelsky2019pwave} shows how strongly a flavor-singlet \(p\)-wave channel can reorganize a bosonic many-body state.  Its number-conserving separable interaction differs from the local quadrature interaction in Eq.~\eqref{eq:three_plus_model_HI}, leaving the fate of its fragmented and pair-coherent regimes on a frustrated lattice as an open question.

Three time reversals act on the annihilation operators as
\begin{equation}
\begin{array}{c|cc|c}
    &\hat b_{\mathbf r}&\hat c_{\mathbf r}&\ii\\
    \hline
    \widehat\Theta_1&\hat b_{\mathbf r}&\hat c_{\mathbf r}&-\ii\\
    \widehat\Theta_2&\ii\hat c_{\mathbf r}&\ii\hat b_{\mathbf r}&-\ii\\
    \widehat\Theta_3&\ii\hat b_{\mathbf r}&-\ii\hat c_{\mathbf r}&-\ii
\end{array}
    \label{eq:three_plus_model_TR}
\end{equation}
Their action on creation operators follows by Hermitian conjugation.  Each one-body map squares to \(+I\), and every pair anticommutes.  The quadratic Hamiltonian is invariant under all three transformations.  Both Hermitian operators inside square brackets in Eq.~\eqref{eq:three_plus_model_HI} are odd under all three time reversals.  Denoting either bracket briefly by \(\hat O\), its Hubbard--Stratonovich decomposition is
\begin{equation}
    \ee^{-\Delta\tau g\hat O^2}
    =\frac{1}{\sqrt{2\pi}}
    \int_{-\infty}^{\infty}d\phi\,\ee^{-\phi^2/2}
    \ee^{\ii\sqrt{2\Delta\tau g}\,\phi\hat O},
    \qquad g=g_1\ \text{or}\ g_2.
    \label{eq:three_plus_HS}
\end{equation}
Antiunitarity reverses \(\ii\), while time reversal reverses \(\hat O\).  Every fixed-field quadratic factor retains the three symmetries in Eq.~\eqref{eq:three_plus_model_TR}.  Equation~\eqref{eq:boson_three_plus_weight} then makes every convergent configuration weight positive.

The most direct weak-coupling reference state is an \(SU(2)\)-symmetric squeezed state with odd-parity singlet-pair correlations and no one-boson condensate.  A particularly interesting competing state places the two flavors near opposite valleys.  At the quadratic instability, if the soft condensate amplitude at \(+\mathbf K\) is a complex two-component spinor \(z\), the singlet pair field selects a partner at \(-\mathbf K\) proportional to \(\ii\sigma_y z^*\).  The two valleys then have equal total populations and opposite flavor polarizations.  Their opposite bond-mass-current patterns cancel, while their bond-flavor-current contributions add.  Coherence between the locked valley amplitudes produces a \(120^\circ\) flavor texture.  This is the paired analogue of the chiral spin superfluid and momentum-space antiferromagnet proposed in Ref.~\cite{liNatuParamekanti2014chiral}.  Opposite projections of its flavor current along a chosen laboratory axis are connected continuously by \(SU(2)\), in accord with the absence of an additional Ising order parameter.

The spinor \(z\) has four real components, and every local \(SU(2)\)-invariant potential for one such spinor depends only on its length.  If its low-energy dynamics is relativistic and no additional soft mode or relevant anisotropy intervenes, a direct continuous transition is therefore a natural candidate for the three-dimensional \(O(4)\) Wilson--Fisher universality class~\cite{pelissettoVicari2002critical}.  Strong fluctuations can instead drive a first-order transition.  They may also produce an intermediate composite phase with a flavor-current vector but no one-boson coherence.  Such a phase would first break \(SU(2)\) down to \(U(1)\), allowing an \(O(3)\)-like onset followed by a second transition that orders the remaining phase.  Current order without single-particle coherence is well established in related frustrated, number-conserving boson models~\cite{dhar2013,zaletel2014prb}, while its flavor-vector counterpart in the present anomalous model remains to be determined.

The exact \(SU(2)\) symmetry also changes the finite-temperature picture.  In two dimensions it precludes conventional long-range flavor order, and the fixed-amplitude spinor manifold \(S^3\) has no ordinary vortices.  Flavor-anisotropic double-valley systems can instead separate a superfluid \(U(1)\) transition from an Ising valley or chirality transition.  Classical Monte Carlo found a narrow chiral nonsuperfluid regime in an orbital lattice model~\cite{liParamekantiHemmerich2014chiral}, whereas a recent shaken-lattice experiment observed the reverse sequence, with valley balance restored before superfluid coherence was lost~\cite{yuWuZhang2025vestigial}.  These contrasting results make the exact \(SU(2)\) line a useful high-symmetry reference point and show how sensitively multivalley melting depends on the allowed order-parameter manifold.

Sign-problem-free simulations can now distinguish these possibilities through valley-resolved one-particle correlations, flavor-current structure factors, pair-quadrature correlations, and the joint probability distribution of the four spinor components.  Correlation-length and Binder-ratio crossings can test \(O(4)\) or split critical scaling, while energy histograms can resolve a weak first-order transition.  The calculation must remain in the region where the bosonic trace converges.  A perfectly balanced semiclassical counterflow can make both positive squares vanish, so these interactions do not by themselves stabilize a mode after the free Bogoliubov kernel becomes negative.  Within the stable region, the model poses a focused nonperturbative question: whether local pair and current interactions turn an \(SU(2)\)-symmetric paired state into a finite-momentum counterflow condensate, generate an intermediate composite phase, or avoid criticality through a first-order transition.

\section{Summary}

Mutually anticommuting time reversals determine the sign of general bosonic and fermionic Gaussian quantum Monte Carlo weights through three physical mechanisms.  The Majorana class \(\{+,-\}\) cancels phases between conjugate sectors.  The Kramers class \(\{-,-\}\) selects one operator from every Nambu-conjugate pair and enforces Kramers degeneracy among them.  The bosonic split-orthogonal class \(\{+,+,+\}\) produces a real propagation that preserves a symmetric matrix with equally many positive and negative directions, which fixes the remaining determinant sign.  This last mechanism has no fermionic counterpart, as demonstrated by the explicit negative two-mode trace in Eq.~\eqref{eq:fermion_three_plus_negative}.  Four time reversals always generate a Majorana pair.

The resulting periodic table applies directly to anomalous Gaussian evolution.  Pairing terms therefore fit naturally within time-reversal-protected auxiliary-field simulations.  The two bosonic models above show how the symmetry algebra can guide Hubbard--Stratonovich design and open sign-problem-free routes to interacting paired phases on frustrated lattices.

\appendix

\section{Equivalence of the fermionic Nambu and Majorana formulations}

For every fermionic mode, introduce two Hermitian Majorana operators,
\begin{equation}
    \hat\gamma_j^1=\hat a_j+\hat a_j^+,
    \qquad
    \hat\gamma_j^2=-\ii(\hat a_j-\hat a_j^+).
    \label{eq:nambu_majorana_operators}
\end{equation}
They obey \(\{\hat\gamma_j^\alpha,\hat\gamma_k^\beta\}=2\delta_{jk}\delta_{\alpha\beta}\).  Collecting first all \(\hat\gamma_j^1\) and then all \(\hat\gamma_j^2\) gives
\begin{equation}
    \hat\gamma=S\hat\chi,
    \qquad
    S=\begin{pmatrix}I_N&I_N\\-\ii I_N&\ii I_N\end{pmatrix},
    \qquad
    \overline S=SC,
    \qquad
    SCS^\trans=2I_{2N}.
    \label{eq:nambu_majorana_basis}
\end{equation}
For the generator and a time reversal \(\Theta=U\mc K\), define their Majorana-basis matrices and the inverse transformation by
\begin{equation}
    A_{\ell,{\rm M}}=SA_\ell S^{-1},
    \qquad
    U_{\rm M}=\overline S\,U S^{-1},
    \qquad
    A_\ell=S^{-1}A_{\ell,{\rm M}}S,
    \qquad
    U=(\overline S)^{-1}U_{\rm M}S.
    \label{eq:majorana_TR_dictionary}
\end{equation}
\begin{samepage}
Anti-linearity accounts for the complex conjugation of \(S\) in \(U_{\rm M}\).  Since \(g_{\rm F}=C\), direct substitution gives the same quadratic generator in Majorana form and maps every condition in Eq.~\eqref{eq:physical_TR_conditions} in both directions:
\begin{align}
    \hat h_\ell
    &=\frac14\hat\gamma^\trans A_{\ell,{\rm M}}\hat\gamma,
    &
    CA_\ell+A_\ell^\trans C=0
    &\quad\Longleftrightarrow\quad
    A_{\ell,{\rm M}}^\trans=-A_{\ell,{\rm M}},
    \label{eq:majorana_generator_dictionary}\\
    C\overline U C=U
    &\quad\Longleftrightarrow\quad
    \overline U_{\rm M}=U_{\rm M},
    &
    U^\trans C U=C
    &\quad\Longleftrightarrow\quad
    U_{\rm M}^\trans U_{\rm M}=I,
    \notag\\
    U^{-1}\overline{A}_\ell U=A_\ell
    &\quad\Longleftrightarrow\quad
    U_{\rm M}^{-1}\overline{A}_{\ell,{\rm M}}U_{\rm M}
    =A_{\ell,{\rm M}}.
    \label{eq:majorana_TR_conditions}
\end{align}
\end{samepage}
Thus a physical Nambu time reversal becomes a real orthogonal Majorana-time-reversal symmetry of every slice.  Using \(\overline U_{\rm M}U_{\rm M}=S(\overline U U)S^{-1}\), its square and its anticommutation relations obey
\begin{align}
    \overline U U=\pm I
    &\quad\Longleftrightarrow\quad
    U_{\rm M}^2=\pm I
    \quad\Longleftrightarrow\quad
    U_{\rm M}^\trans=\pm U_{\rm M},
    \label{eq:majorana_square_dictionary}\\
    \overline U_bU_a+\overline U_aU_b=0
    &\quad\Longleftrightarrow\quad
    U_{{\rm M},b}U_{{\rm M},a}
    +U_{{\rm M},a}U_{{\rm M},b}=0.
    \label{eq:majorana_anticommutation_dictionary}
\end{align}
Real symmetric and real antisymmetric orthogonal matrices are precisely the positive- and negative-square Majorana-time-reversal matrices of Ref.~\cite{li2016prl}.  Applying the inverse formulas in Eq.~\eqref{eq:majorana_TR_dictionary} maps the right-hand conditions back to the fermionic Nambu conditions on the left.  Hence the fermionic \(\{+,-\}\) class maps one-to-one onto the Majorana class of Ref.~\cite{li2016prl}.  The fermionic \(\{-,-\}\) class maps one-to-one onto its Kramers class.

\bibliography{references}

\end{document}
